% Source: https://artofproblemsolving.com/community/c3h2578013_flipping_a_coin
% Stable identity: aops-topic-2578013-q01
Each pair (coin, $d$) has probability $\dfrac12\cdot\dfrac1{10}=\dfrac1{20}$.

\emph{``Multiple''.} The choices are $d,2d,\dots,9d$, nine in all. Since $d\ge2012$, the only way to get $2020$ is $d=2020$ with multiple $1\cdot d$, and $2$ is impossible. Contribution: $\dfrac1{20}\cdot\dfrac19$.

\emph{``Divisor''.} The number $2$ is a divisor exactly for even $d$, and $2020$ is a divisor only for $d=2020$. The relevant divisor counts are
\[\begin{array}{c|ccccc}
d&2012=2^2\cdot503&2014=2\cdot19\cdot53&2016=2^5\cdot3^2\cdot7&2018=2\cdot1009&2020=2^2\cdot5\cdot101\\\hline
\tau(d)&6&8&36&4&12
\end{array}\]
For $d=2020$ both $2$ and $2020$ are favourable. Contribution:
\[\dfrac1{20}\left(\dfrac16+\dfrac18+\dfrac1{36}+\dfrac14+\dfrac2{12}\right).\]

In total,
\[\dfrac1{20}\left(\dfrac19+\dfrac16+\dfrac18+\dfrac1{36}+\dfrac14+\dfrac16\right)=\dfrac1{20}\cdot\dfrac{8+12+9+2+18+12}{72}=\dfrac{1}{20}\cdot\dfrac{61}{72}=\boxed{\dfrac{61}{1440}}.\]
